% Generated draft: requires mathematical and source review.
\documentclass[11pt]{article}
\usepackage[margin=25mm]{geometry}
\usepackage{fontspec}
\setmainfont{Latin Modern Roman}
\usepackage{amsmath,amssymb,mathtools,unicode-math}
\setmathfont{Latin Modern Math}
\usepackage{enumitem,needspace,xurl,hyperref}
\setlength{\parindent}{0pt}
\setlength{\parskip}{5pt}
\setlength{\emergencystretch}{3em}
\newcommand{\N}{\mathbb{N}}
\newcommand{\Z}{\mathbb{Z}}
\newcommand{\Q}{\mathbb{Q}}
\newcommand{\R}{\mathbb{R}}
\newcommand{\C}{\mathbb{C}}
\newcommand{\F}{\mathbb{F}}
\newcommand{\A}{\mathbb{A}}
\newcommand{\PP}{\mathbb{P}}
\newcommand{\E}{\mathbb{E}}
\newcommand{\eps}{\varepsilon}
\newcommand{\dd}{\,\mathrm{d}}
\begin{document}
\section*{UnsolvedMath ID: 30004156}
\textbf{Source problem code:} catalogue entry
\subsection{Definitions and mathematical statement}
Use the standard convention that a line in an equiangular line system is a one-dimensional linear subspace of \(\R^n\), hence passes through the origin. For distinct such lines \(\ell,m\), choose unit vectors \(u\in\ell\) and \(v\in m\), and define their unoriented angle by
\[
\angle(\ell,m)=\arccos\bigl(|\langle u,v\rangle|\bigr)\in(0,\pi/2].
\]
The absolute value makes this independent of the choices of unit vectors. A finite set \(\mathcal L\) of distinct lines is equiangular with common angle \(\theta\) if
\[
\angle(\ell,m)=\theta
\qquad\text{for every distinct }\ell,m\in\mathcal L.
\]
For each positive integer \(n\) and each prescribed \(\theta\in(0,\pi/2]\), define
\[
M(n,\theta)=\max\left\{|\mathcal L|:\mathcal L\text{ is an equiangular set of lines in }\R^n\text{ with common angle }\theta\right\}.
\]
Sets with at most one line satisfy the pairwise condition vacuously, so they are allowed even when no pair of distinct lines realizes the prescribed angle. Define also
\[
M(n)=\max_{\theta\in(0,\pi/2]}M(n,\theta).
\]
The problem is to determine \(M(n)\) for every positive integer \(n\), and \(M(n,\theta)\) for every positive integer \(n\) and every prescribed \(\theta\in(0,\pi/2]\).

The supplied record does not explicitly specify whether lines pass through the origin, the angle convention, or endpoint restrictions. The formulation above uses the conventional interpretation of equiangular line systems, includes the orthogonal endpoint \(\theta=\pi/2\), and excludes the degenerate angle \(0\), which cannot occur between distinct linear lines. These conventions require confirmation against the cited report.
\subsection{Short English statement}
Determine the largest possible number of distinct equiangular lines through the origin in \(\R^n\), both when the common unoriented angle may vary and when it is prescribed. The source labels the general problem open; this is a supplied status, not an independently verified assessment.
\subsection{Sources}
\begin{itemize}
\item Supplied local UnsolvedMath registry record, permanent ID \(30004156\), title \emph{Maximum Equiangular Line Systems in Euclidean Space}, source problem code \texttt{catalogue entry}, in the Oberwolfach Reports Open Problems collection. No individual proposer is supplied; the recorded proposed year is 2019.
\item Supplied report citation: \emph{Mixed-integer Nonlinear Optimization: a hatchery for modern mathematics} (2020), Oberwolfach Reports, pp.\ 1573--1637, DOI \texttt{10.4171/owr/2019/26}. Source URL and literature-source label ``OWR primary report'': \url{https://doi.org/10.4171/owr/2019/26}.
\item Supplied literature source, labeled ``equiangular lines in low dimensions'': \url{https://arxiv.org/abs/2002.08085}. The supplied record does not give authors or a more detailed bibliographic citation.
\item Literature assessment dated exactly \texttt{2026-08-22}: the supplied curation labels the problem open and states that maximum sizes are known in many dimensions, strong general bounds exist, and exact values remain unknown in general, including unresolved dimensions documented in the cited primary work. This is supplied dated curation derived from the OWR record, not an independently verified literature assessment; time-sensitive status and cited claims require independent verification.
\item The registry labels the statement \texttt{corrected\_verified} and says that mathematical content was preserved while notation, typography, and sentence structure were repaired for publication. This supplied editorial label does not establish independent mathematical or source verification of this draft.
\end{itemize}
\end{document}
